
\chapter{Datasets and Preprocessing}
\label{ch:datasets}

\section{PlantVillage}

PlantVillage \cite{hughes2015} supplies 38 classes covering 14 crop species with
healthy and diseased variants. Images were collected to isolate the leaf against a
uniform background under controlled lighting, which is what makes the benchmark clean
and what makes its assumptions narrow.

\begin{table}[H]
\centering
\caption{PlantVillage splits used in this study}
\label{tab:pvsplits}
\tabfont
\begin{tabularx}{\textwidth}{@{}lrrY@{}}
\toprule
\textbf{Split} & \textbf{Images} & \textbf{Classes} & \textbf{Role} \\
\midrule
Train & 36,762 & 38 & fitting \\
Validation & 6,544 & 38 & checkpoint selection and calibration \\
Development total & 43,306 & 38 & train plus validation \\
Locked test & 10,709 & 38 & final evaluation, scored once \\
\bottomrule
\end{tabularx}
\end{table}

Splits were constructed at the group level so that no near-duplicate occupies two
sides of a split. The frozen split definition is stored in
`data/manifests/split\_freeze.json` and is checked by the test suite; a change to the
split breaks the tests rather than silently shifting results.

\section{PlantDoc}

PlantDoc \cite{singh2020} represents field conditions: cluttered backgrounds,
variable illumination, leaves partially occluded by other leaves, and photographs
taken at working distance rather than for a camera.

\begin{table}[H]
\centering
\caption{PlantDoc as used for field evaluation}
\label{tab:pdsplits}
\tabfont
\begin{tabularx}{\textwidth}{@{}lrrY@{}}
\toprule
\textbf{Property} & \textbf{Value} & & \textbf{Note} \\
\midrule
Dataset classes & 28 & & all assigned a PlantVillage target \\
Test images & 236 & & the field evaluation set \\
Classes with test images & 27 & & one class has none in this split \\
Validation split & none & & so field data is never used for selection \\
\bottomrule
\end{tabularx}
\end{table}

\section{Label mapping}

The two datasets do not share a label space, so every PlantDoc class was assigned a
PlantVillage target explicitly, and the assignment is stored with its justification.

\begin{table}[H]
\centering
\caption{PlantDoc to PlantVillage mapping decisions}
\label{tab:mapping}
\tabfont
\begin{tabularx}{\textwidth}{@{}lY@{}}
\toprule
\textbf{Decision} & \textbf{Count} \\
\midrule
Mapped, status confirmed & 25 \\
Mapped, confirmed with caveat & 3 \\
Unmapped and excluded from scoring & 0 \\
\midrule
Relationship labelled ``same disease'' & 12 \\
Relationship labelled ``healthy class'' & 10 \\
Relationship labelled ``synonym'' & 4 \\
Relationship labelled ``likely synonym'' & 1 \\
Relationship labelled ``assumed synonym'' & 1 \\
\bottomrule
\end{tabularx}
\end{table}

Three mappings carry a caveat rather than a confirmation. They are retained because
the alternative, dropping the class, discards real test images and biases the field
estimate towards classes that happen to be easy to map. The caveat is recorded in
`data/manifests/plantdoc\_to\_plantvillage\_map.csv` beside the mapping itself.

An earlier version of this mapping covered only part of the label space and produced
a meaningless macro-F1, because predicting a class outside the scored subset was
counted as a miss against a class that was not being scored. That pathology is
documented in Chapter~\ref{ch:limitations} and the mapping was completed before any
number in this report was produced.

\section{Preprocessing}

Images are resized to the model input with bilinear interpolation, converted to
tensors, and normalised with the ImageNet mean and standard deviation stored inside
the checkpoint. The normalisation constants travel with the weights, so an inference
path cannot silently disagree with the training path.

Training-only augmentation is described in Table~\ref{tab:protocol}. No augmentation
is applied at evaluation, on either dataset.

\section{Dataset limitations}

PlantVillage is not a sample of field conditions and was not collected to be one.
Results on it describe behaviour under the conditions it encodes.

PlantDoc is small. At 236 test images the 95\,\% interval on a single accuracy
estimate is roughly ten percentage points wide, which is why paired comparisons and
replication are used rather than a single number. It is also not a random sample of
Indian field conditions; it is the best available proxy.

Neither dataset supports a claim about severity, multiple concurrent diseases, or
disease on species absent from the label set.
